\scanpage{146}
\deleted{Alors $\mathrm{SO}(q)$ est fibré sur}

\deleted{$[$Considérons}
\[
\begin{aligned}
&\deleted{u\longmapsto u(F)}\\
&\deleted{\mathrm{SO}(q)\longrightarrow\underline G_{\{n\}}}
\end{aligned}
\]
\deleted{il suffit de prouver que ce morphisme est lisse à fibres géom $0$-conn. Plus gén., On va prouver que pour tt}
\deleted{\uline{Lissité}}
\[
\deleted{D\in\underline G_H(E)}\qquad
\left\{\begin{aligned}
&\deleted{\mathrm{SO}(q)\longrightarrow\underline G_H}\\
&\deleted{u\longmapsto u(D)}
\end{aligned}\right.
\]
\deleted{est lisse à fibres géom $0$-conn. (donc $\underline G_H$ espace homogène sous $\mathrm{SO}(q)$). Pour transitivité, \unsure{SPS|OPS} $H=[1,n]$}

\deleted{Comme $\mu_2$ opère \unsure{trivialmt|transitivement} sur les $\underline G_H$, $\underline{\mathrm O}(q)$ opère via $\underline{\mathrm{SO}}(q)$}
\deleted{\uline{Lissité} \unsure{groupe des translations}. Posons}
\[
\deleted{D=\{E_1\subset\ldots\subset E_n\subset\underbrace{E}_{\unsure{S}}\}}\qquad\deleted{\mathrm{rg}\,E_i=1.}
\]
\deleted{\unsure{Suggestion|Suggested}} \deleted{\unsure{q}} \deleted{Soit}

\deleted{On voit aisément \inserted{(cf. lemme plus bas)} qu'il est associé à une base standard, $e_0,e_1\ldots e_n,e_1^*\ldots e_n^*,$ \deleted{$e_0$} avec $E_i$ engendré par $e_1\ldots e_i$. \unsure{Dé|A}}

\deleted{\unsure{À|Si} un $D'=(E'_1\subset\ldots\subset E'_n)\in G_{[1,n]}$ est associé : $e'_0,e'_1\ldots e'_n\ e_1^{\prime *}\ldots e_n^{\prime *}$. Il y a un autom. (unique) \inserted{de} $(E,q)$ $e_i\longrightarrow e'_i$, $e_i^*\longrightarrow e_i^{\prime *}$, Il donne}
% archive sheet number: 145
